\section*{38. Avec R. Brauer et H. Hasse: preuve d'un théorème principal dans la théorie des algèbres}

\begin{center}
\emph{Journal f. d. reine u. angew. Math. 167 (1932), p. 399--404}
\end{center}

\begin{center}
\textbf{Preuve d'un théorème principal dans la théorie des algèbres.}\\[0.5em]
Par \emph{R. Brauer} à Koenigsberg, \emph{H. Hasse} à Marbourg et \emph{E. Noether} à Goettingen\footnote{La rédaction de cette note fut prise en charge par H. Hasse.}.
\end{center}

Enfin, nos efforts conjoints ont permis d'établir le théorème suivant, fondamental pour la théorie structurale des algèbres sur les corps de nombres algébriques et au-delà:

\paragraph{Théorème principal.}
\emph{Toute algèbre à division normale sur un corps de nombres algébrique est cyclique (ou, comme on dit aussi, du type de Dickson).}

Nous avons le plaisir particulier de dédier ce résultat --- succès dû pour l'essentiel à la méthode \(p\)-adique --- à M. Kurt Hensel, fondateur de cette méthode, à l'occasion de son soixante-dixième anniversaire.

Notre preuve comporte trois réductions, dont chacun de nous a fourni une.\footnote{Elles sont reproduites dans l'ordre de leur genèse, qui est inverse de l'ordre systématique.}

\paragraph{1.}
La première réduction fut donnée par H. Hasse sur la base de la théorie des algèbres cycliques sur les corps de nombres algébriques qu'il avait récemment développée.\footnote{H. Hasse, \emph{Theorie der zyklischen Algebren ueber einem algebraischen Zahlkoerper}, Goett. Nachr. 1931. Un exposé détaillé des preuves, qui développe notamment la théorie des corps de décomposition et des produits croisés élaborée par E. Noether dans un cours, fondamentale là comme ici, paraîtra prochainement dans les Trans. Amer. Math. Soc. sous le titre \emph{Theory of cyclic algebras over an algebraic number field}. Ce travail sera cité dans la suite par H. H; 1--6 constituent, à la différence de la langue près, la première note.}

\paragraph{Réduction 1.}
Il suffit, pour démontrer le théorème principal, d'établir que:
\[
\text{I. Toute algèbre sur }\Omega\text{ qui se décompose partout est } \sim \Omega .
\]
Pour abréger l'expression, ici comme dans la suite, ``algèbre \(A\) sur \(\Omega\)'' signifiera toujours ``algèbre simple normale \(A\) sur le corps de nombres algébrique \(\Omega\)'', c'est-à-dire un système hypercomplexe simple de centre \(\Omega\). De plus, \(A\sim\Omega\) signifie que \(A\) est une algèbre complète de matrices sur \(\Omega\), donc l'appartenance de \(A\) à la classe spéciale, au sens de la similitude -- c'est-à-dire l'égalité des algèbres à division associées sur \(\Omega\) --, déterminée par \(\Omega\) lui-même comme algèbre à division sur \(\Omega\) [H, 5]. Enfin, par algèbre sur \(\Omega\) se décomposant partout, on entend une algèbre telle que, pour toute place première \(\mathfrak p\) de \(\Omega\), l'extension \(\mathfrak p\)-adique, c'est-à-dire l'algèbre obtenue par extension du corps des coefficients \(\Omega\) au corps \(\mathfrak p\)-adique correspondant \(\Omega_{\mathfrak p}\), se décompose:
\[
A_{\mathfrak p}\sim \Omega_{\mathfrak p}.
\]

\paragraph{Preuve.}
Soit \(D\) une algèbre à division sur \(\Omega\). Si \(Z\) est alors un corps cyclique sur \(\Omega\) tel que, pour toute place première \(\mathfrak p\) de \(\Omega\), le \(\mathfrak p\)-degré de \(Z\) soit un multiple du \(\mathfrak p\)-indice de \(D\) [cf. la note relative à H, 17 Bb], alors, pour les diviseurs premiers correspondants \(\mathfrak P\) dans \(Z\), le corps \(\mathfrak P\)-adique \(Z_{\mathfrak P}\) est chaque fois corps de décomposition de \(D_{\mathfrak p}\) [H, 18.1]. Il en résulte que l'algèbre \(D_Z\), obtenue par extension du corps des coefficients de \(D\) à \(Z\), se décompose partout sur \(Z\). Si I est déjà acquis, pour tout \(\Omega\), il s'ensuit donc que \(D_Z\sim Z\). Cela signifie que \(D\) possède le corps de décomposition cyclique \(Z\), donc qu'elle est représentable cycliquement [H, 5]. Alors \(D\) possède aussi un tel corps de décomposition cyclique dont le degré coïncide avec le degré de \(D\) lui-même; autrement dit, \(D\) est cyclique [H, 6, théorème 6]. Sous l'hypothèse I, le théorème principal est donc démontré.

\paragraph{2.}
Pour la deuxième réduction, R. Brauer donna l'impulsion décisive en communiquant par lettre à H. Hasse comment la question voisine de la valeur exacte de l'exposant d'une algèbre à division (question qui est d'ailleurs résolue en même temps par le théorème principal, voir ci-dessous) peut être ramenée, au moyen du théorème des groupes de Sylow, au cas d'un corps de décomposition résoluble. Cette idée put alors être utilisée aussi pour une réduction correspondante de la question de cyclicité.

\paragraph{Réduction 2.}
L'assertion I est démontrée s'il est montré que:
\[
\text{II. Toute algèbre partout décomposée ayant un corps de décomposition résoluble sur }\Omega\text{ est }\sim\Omega .
\]

\paragraph{Preuve.}
Soit \(A\) une algèbre partout décomposée sur \(\Omega\). Si \(K\) est un corps de décomposition galoisien pour \(A\), si de plus \(p\) est un nombre premier quelconque et \(\Sigma\) l'un des corps de Sylow de \(K\) sur \(\Omega\) relatifs à \(p\) (c'est-à-dire le corps des invariants d'un groupe de Sylow relatif à \(p\) du groupe de Galois), alors \(A_{\Sigma}\) est une algèbre également partout décomposée sur \(\Sigma\), qui possède le corps de décomposition résoluble \(K\). Si II est déjà acquis (pour tout \(\Omega\)), il s'ensuit donc que \(A_{\Sigma}\sim \Sigma\). Cela signifie que \(A\) possède le corps de décomposition \(\Sigma\), de degré premier à \(p\). L'indice de \(A\), étant diviseur de ce degré, est donc premier à \(p\). Comme cela vaut pour tout nombre premier \(p\), il est donc égal à 1, c'est-à-dire \(A\sim\Omega\). Sous l'hypothèse II, l'assertion I est donc démontrée.\footnote{L'idée de la réduction à un corps de décomposition résoluble au moyen du théorème des groupes de Sylow avait déjà été employée auparavant par R. Brauer, pour montrer que tout diviseur premier de l'indice apparaît aussi dans l'exposant (Ueber den Zusammenhang von arithmetischen und invariantentheoretischen Eigenschaften von Gruppen linearer Substitutionen, Berl. Akad.-Ber. 1926). Plus récemment A. A. Albert a développé, pour cette idée et plus généralement pour une série de théorèmes généraux de la théorie de R. Brauer et E. Noether, des preuves simples indépendantes de la théorie des représentations (1. \emph{On direct products, cyclic division algebras, and pure Riemann matrices}; 2. \emph{On direct products}; tous deux dans Trans. Amer. Math. Soc. 33 (1931); pour la réduction ici en question, voir surtout le théorème 23 dans 2.).

\emph{Addition pendant l'impression.} En outre A. A. Albert, ayant en sa possession la communication épistolaire de H. Hasse selon laquelle le théorème principal avait été démontré par lui pour les algèbres abéliennes (voir plus bas dans le texte), en a immédiatement déduit, indépendamment de nous, les faits suivants:

a) le théorème principal pour les degrés de la forme \(2^e\),

b) le théorème 1 ci-dessous (exposant = indice),

c) outre l'idée fondamentale de la réduction 2, celle aussi de la réduction 3 suivante, naturellement sans relation avec la réduction 1, et par conséquent avec le résultat: pour les algèbres à division \(D\) de degré \(p^e\), puissance d'un nombre premier, sur \(\Omega\), il existe un corps d'extension \(\Omega'\) de degré premier à \(p\) sur \(\Omega\), tel que \(D_{\Omega'}\) soit cyclique.

Ces trois résultats sont naturellement dépassés par notre démonstration désormais achevée du théorème principal. Ils montrent cependant qu'A. A. Albert a lui aussi contribué indépendamment à sa preuve. Enfin, après avoir pris connaissance de notre démonstration, A. A. Albert a remarqué que notre théorème central I découle en quelques lignes des théorèmes 13, 10 et 9 d'un travail de lui en cours d'impression (Bull. Amer. Math. Soc. 37 (1931)). La preuve de ces théorèmes repose essentiellement sur les mêmes raisonnements que nos réductions 2 et 3.}

\paragraph{3.}
La troisième réduction --- dont H. Hasse reconnut, une fois en possession des deux premières, qu'elle conduisait à la preuve définitive --- fut donnée par E. Noether à la suite d'une communication de H. Hasse. Celle-ci démontrait le théorème principal dans le cas particulier d'un corps de décomposition abélien, au moyen de la réduction générale 1 et d'une réduction explicite du système de facteurs correspondant; E. Noether en dégagea précisément le noyau sous la forme de la troisième réduction. Indépendamment d'E. Noether, R. Brauer avait d'ailleurs déjà envisagé cette troisième réduction.

\paragraph{Réduction 3.}
L'assertion II est démontrée s'il est montré que:
\[
\begin{gathered}\text{III. Toute algèbre partout décomposée ayant un corps de décomposition cyclique}\\\text{de degré premier sur }\Omega\text{ est }\sim\Omega .\end{gathered}
\]

\paragraph{Preuve.}
Soit \(A\) une algèbre partout décomposée ayant un corps de décomposition résoluble \(K\) sur \(\Omega\). Soit
\[
K=\Lambda_0>\Lambda_1>\cdots>\Lambda_r=\Omega
\]
une chaîne de corps entre \(K\) et \(\Omega\), telle que \(\Lambda_i\) soit toujours cyclique de degré premier sur \(\Lambda_{i+1}\). Alors \(A_{\Lambda_1}\) est d'abord une algèbre également partout décomposée sur \(\Lambda_1\), qui possède le corps de décomposition cyclique de degré premier \(K=\Lambda_0\). Si III est déjà acquis (pour tout \(\Omega\)), il s'ensuit donc que \(A_{\Lambda_1}\sim \Lambda_1\). Cela signifie que \(\Lambda_1\) est corps de décomposition pour \(A\). On démontre maintenant de la même façon, en partant de \(\Lambda_1\) au lieu de \(\Lambda_0\), que \(A_{\Lambda_2}\sim\Lambda_2\), et ainsi de suite, jusqu'à obtenir finalement \(A_{\Lambda_r}\sim\Lambda_r\), c'est-à-dire \(A\sim\Omega\). Sous l'hypothèse III, l'assertion II est donc démontrée.

\paragraph{4.}
Or III est acquis d'après les résultats de H. Hasse [H, 24]. Par conséquent, en remontant les réductions 3, 2, 1, on obtient l'exactitude du théorème principal.

E. Noether indique encore à ce propos ce qui suit: l'exactitude de III, plus généralement pour des corps de décomposition cycliques de degré quelconque, revient au théorème de norme de Hasse\footnote{\label{fn:p38-hasse-normensatz}Cité dans H, 3.11; démontré dans: H. Hasse, \emph{Beweis eines Satzes und Widerlegung einer Vermutung ueber das allgemeine Normenrestsymbol}, Goett. Nachr. 1931.}. Pour la réduction 3, seul le cas spécial du degré premier est cependant nécessaire, donc le théorème de norme originel de Hilbert-Furtwaengler\footnote{Voir H. Hasse, Bericht ueber neuere Untersuchungen und Probleme in der Theorie der algebraischen Zahlkoerper II, Jahresber. der D. M.-V., Erg.-Bd. 6 (1930), \S 8, ainsi que la littérature qui y est citée.}. Ainsi la réduction III fournit une nouvelle preuve simple du théorème de norme de Hasse. Cette preuve se déroulerait ainsi:

Soit \(Z\) un corps cyclique sur \(\Omega\) et \(\alpha\) un nombre de \(\Omega\), pour lequel le symbole de reste normique
\[
\left(\frac{\alpha,Z}{\mathfrak p}\right)=1
\]
vaut pour toute place première \(\mathfrak p\) de \(\Omega\). Alors toute algèbre cyclique
\[
A=(\alpha,Z)
\]
\footnote{Dans la notation de H, 1. -- L'indication de l'automorphisme \(S\) de \(Z\) lié à \(\alpha\) a été omise ici comme inessentielle.} se décompose partout [H, 17.7]. D'après la réduction 3 il s'ensuit donc, en utilisant seulement le théorème de norme de Hilbert-Furtwaengler, que \(A\sim\Omega\). Mais alors \(\alpha\) est la norme d'un élément de \(Z\) [H, 15.4].

En rapport avec le théorème de norme, nous remarquons encore que l'assertion I doit être regardée comme sa véritable généralisation aux cas supérieurs (même non abéliens), tandis que la généralisation littérale, comme H. Hasse l'a montré\textsuperscript{\ref{fn:p38-hasse-normensatz}}, n'est pas vraie en général.

\begin{center}
\textbf{Conséquences (H. Hasse).}
\end{center}

\paragraph{5.}
Comme conséquence la plus immédiate du théorème principal, signalons que la théorie des algèbres cycliques sur les corps de nombres algébriques développée par H. Hasse a désormais acquis la portée d'une théorie générale de la structure et des invariants des algèbres simples normales (en particulier des algèbres à division normales) sur les corps de nombres algébriques. En particulier, le résultat suivant est maintenant démontré en toute généralité:

\paragraph{Théorème 1.}
L'exposant d'une algèbre simple normale sur un corps de nombres algébrique est égal à son indice.

De plus, l'assertion I donne le fait remarquable suivant:

\paragraph{Théorème 2.}
L'idéal fondamental d'une vraie algèbre à division normale sur \(\Omega\) est divisible par au moins une place première de \(\Omega\), et même par au moins deux.

Ici l'idéal fondamental est défini par exemple comme la norme (réduite) de la différente (réduite).\footnote{Dans le cas spécial des algèbres de quaternions rationnelles, cela revient à la notion de ``nombre fondamental'' introduite par H. Brandt (Idealtheorie in Quaternionenalgebren, Math. Ann. 99 (1928)). Comme il ne s'agit pas ici d'un nombre, mais d'un idéal, il fallait dire ``idéal fondamental''; on ne doit pas confondre cela avec les dénominations de Dedekind idéal fondamental et nombre fondamental pour la différente et le discriminant, qui justement pour la raison indiquée se révèlent inutilisables pour les corps relatifs. -- Pour la définition de la différente (réduite), voir la note suivante. La norme (réduite) d'un idéal est également définie par E. Noether ``place première par place première'', plus précisément, conformément au fait que, pour les places premières individuelles, tout idéal est principal, tout simplement par formation des normes de nombres (réduites) des nombres de base d'idéaux principaux pour les places premières individuelles.} En outre, une place première infinie est incluse dans l'idéal fondamental si et seulement si, pour elle, l'algèbre se réduit à l'algèbre de quaternions (et non au corps des nombres réels ou complexes).

\paragraph{Preuve.}
D'après les résultats de H. Hasse\footnote{\label{fn:p38-hasse-padisch}H. Hasse, \emph{Ueber \(p\)-adische Schiefkoerper und ihre Bedeutung fuer die Arithmetik hyperkomplexer Zahlsysteme}, Math. Ann. 104 (1931); voir là les théorèmes 42, 59.}, une place première \(\mathfrak p\) de \(\Omega\) divise l'idéal fondamental si et seulement si le \(\mathfrak p\)-indice est différent de \(1\). Si tous les \(\mathfrak p\)-indices étaient égaux à \(1\), l'algèbre se décomposerait partout, et d'après l'assertion I il n'y aurait alors pas de vraie algèbre à division. (Le fait qu'il ne puisse pas non plus y avoir un seul \(\mathfrak p\)-indice différent de \(1\) résulte de ce que les symboles de reste normique, dont les ordres sont les \(\mathfrak p\)-indices [H, 17.7], satisfont au théorème du produit (loi de réciprocité) [H, 3.8].)

\paragraph{6.}
En outre, l'assertion I permet de déterminer (caractérisation arithmétique) tous les corps de décomposition \(K\) appartenant à une algèbre \(A\) sur \(\Omega\). Dans une généralisation exacte de l'état de choses déjà établi par H. Hasse dans le cas spécial où \(A,K\) sont cycliques [H, 6, théorème 2], on obtient:

\paragraph{Théorème 3.}
Pour une algèbre \(A\) sur \(\Omega\), un corps de nombres algébrique \(K\) sur \(\Omega\) est corps de décomposition si et seulement si, pour tous les diviseurs premiers \(\mathfrak P_i\) dans \(K\) de toutes les places premières \(\mathfrak p\) de \(\Omega\), le \(\mathfrak P_i\)-degré \(n_{\mathfrak P_i}\) de \(K\) est chaque fois un multiple du \(\mathfrak p\)-indice \(m_{\mathfrak p}\) de \(A\).

Ici le \(\mathfrak P_i\)-degré de \(K\) est le degré du corps \(\mathfrak P_i\)-adique \(K_{\mathfrak P_i}\), correspondant à \(\mathfrak P_i\), sur le corps \(\mathfrak p\)-adique \(\Omega_{\mathfrak p}\), c'est-à-dire, si
\[
\mathfrak p=\prod_i \mathfrak P_i^{\,e_{\mathfrak P_i}},
        \qquad N_{K/\Omega}(\mathfrak P_i)=\mathfrak p^{\,f_{\mathfrak P_i}},
\]
est la décomposition de \(\mathfrak p\) dans \(K\), le produit du degré \(f_{\mathfrak P_i}\) par l'ordre de ramification \(e_{\mathfrak P_i}\) de \(\mathfrak P_i\) relativement à \(\mathfrak p\), \(n_{\mathfrak P_i}=f_{\mathfrak P_i}e_{\mathfrak P_i}\) [H, 4]. Et le \(\mathfrak p\)-indice de \(A\) est l'indice \(m_{\mathfrak p}\) de l'extension \(\mathfrak p\)-adique \(A_{\mathfrak p}\) [H, 6].

\paragraph{Preuve.}
Dire que \(K\) est corps de décomposition pour \(A\), c'est-à-dire \(A_K\sim K\), équivaut d'après l'assertion I à dire que \(A_{K_{\mathfrak P_i}}\sim K_{\mathfrak P_i}\) pour tout \(\mathfrak P_i\).

Pour une place première finie \(\mathfrak p\), soit
\[
        A_{\mathfrak p}\sim(\pi,W_{\mathfrak p})
\]
la représentation cyclique arithmétiquement distinguée de \(A_{\mathfrak p}\) [H, 16; en outre loc. cit.\textsuperscript{\ref{fn:p38-hasse-padisch}}, théorème 38], donc \(W_{\mathfrak p}\) le corps non ramifié de degré \(m_{\mathfrak p}\) sur \(\Omega_{\mathfrak p}\) et \(\pi\) un nombre de \(\Omega_{\mathfrak p}\) divisible exactement par \(\mathfrak p^1\). Soit de plus \(K_{\mathfrak P_i}W_{\mathfrak p}\) le composé et
\[
        \Delta_{\mathfrak P_i}=K_{\mathfrak P_i}\cap W_{\mathfrak p}
\]
l'intersection de \(W_{\mathfrak p}\) et \(K_{\mathfrak P_i}\). Comme le sous-corps non ramifié maximal contenu dans \(K_{\mathfrak P_i}\) est de degré \(f_{\mathfrak P_i}\) sur \(\Omega_{\mathfrak p}\), et comme il n'existe, pour chaque degré, qu'un seul corps non ramifié sur \(\Omega_{\mathfrak p}\), \(\Delta_{\mathfrak P_i}\) est
\[
        \text{de degré }d_{\mathfrak P_i}=(m_{\mathfrak p},f_{\mathfrak P_i})\text{ sur }\Omega_{\mathfrak p}.
\]
Par conséquent \(K_{\mathfrak P_i}W_{\mathfrak p}\) est de degré \(m_{\mathfrak p}/d_{\mathfrak P_i}\) sur \(K_{\mathfrak P_i}\), et non ramifié.

Or
\[
        A_{K_{\mathfrak P_i}}=(A_{\mathfrak p})_{K_{\mathfrak P_i}}\sim(\pi,K_{\mathfrak P_i}W_{\mathfrak p})
\]
[H, 15.5]. Par suite, \(A_{K_{\mathfrak P_i}}\sim K_{\mathfrak P_i}\) équivaut à ce que \(\pi\) soit norme de \(K_{\mathfrak P_i}W_{\mathfrak p}\) relativement à \(K_{\mathfrak P_i}\). Mais, en raison du caractère non ramifié de \(K_{\mathfrak P_i}W_{\mathfrak p}\) sur \(K_{\mathfrak P_i}\), cela est le cas si et seulement si l'ordre \(e_{\mathfrak P_i}\) de \(\pi\) en \(\mathfrak P_i\) est un multiple du degré \(m_{\mathfrak p}/d_{\mathfrak P_i}\) de \(K_{\mathfrak P_i}W_{\mathfrak p}\) sur \(K_{\mathfrak P_i}\), ou encore, ce qui revient au même à cause de
\[
\left(\frac{m_{\mathfrak p}}{d_{\mathfrak P_i}},\frac{f_{\mathfrak P_i}}{d_{\mathfrak P_i}}\right)=1,
\]
si
\[
\frac{f_{\mathfrak P_i}}{d_{\mathfrak P_i}}e_{\mathfrak P_i}=\frac{n_{\mathfrak P_i}}{d_{\mathfrak P_i}}
\]
est un multiple de \(m_{\mathfrak p}/d_{\mathfrak P_i}\), c'est-à-dire si \(n_{\mathfrak P_i}\) est un multiple de \(m_{\mathfrak p}\), comme il était affirmé.

Dans le cas où \(\mathfrak p\) est une place première infinie et où l'on n'est pas dans le cas trivial \(m_{\mathfrak p}=1\), on a \(m_{\mathfrak p}=2\) et \(A_{\mathfrak p}\) est l'algèbre de quaternions sur le corps réel \(\Omega_{\mathfrak p}\). Dire que \(A_{K_{\mathfrak P_i}}=(A_{\mathfrak p})_{K_{\mathfrak P_i}}\sim K_{\mathfrak P_i}\) équivaut alors à dire que \(K_{\mathfrak P_i}\) est le corps des nombres complexes, c'est-à-dire à \(n_{\mathfrak P_i}=f_{\mathfrak P_i}=2=m_{\mathfrak p}\) (\(e_{\mathfrak P_i}\) n'intervient pas ici), ce qui donne ici encore l'assertion (\(n_{\mathfrak P_i}\), comme \(m_{\mathfrak p}\), ne peut prendre que les valeurs 1 ou 2).

\paragraph{7.}
On arrive à des résultats significatifs si l'on inverse la question, résolue dans le théorème 3, de tous les corps de décomposition \(K\) d'une algèbre fixe \(A\): pour un corps algébrique fixe \(K\) sur \(\Omega\), on considère l'ensemble de toutes les algèbres \(A\) sur \(\Omega\) décomposées par \(K\). Ces algèbres \(A\) forment un sous-groupe \(\mathfrak K\), déterminé par \(K\), du groupe de R. Brauer \(\mathfrak A\) de toutes les algèbres (plus exactement de toutes les classes d'algèbres semblables) sur \(\Omega\).\footnote{R. Brauer, \emph{Ueber Systeme hyperkomplexer Groessen}, Jahresber. d. D. M.-V. 38 (1929), p. 47/48. -- Voir aussi H, 13.1.} Si l'on suppose \(K\) galoisien sur \(\Omega\), on parvient ainsi à des théorèmes que l'on peut regarder comme des généralisations de théorèmes principaux de la théorie des corps de classes (théorie des corps de nombres relativement abéliens) aux corps de nombres relativement galoisiens généraux:

\paragraph{Théorème 4.}
\emph{(Théorème de décomposition).} Le degré relatif \(f\) des diviseurs premiers dans \(K\) d'un idéal premier \(\mathfrak p\) de \(\Omega\) ne divisant pas le discriminant relatif de \(K\) est égal au plus petit exposant pour lequel
\[
        A_{\mathfrak p}^{\,f}\sim\Omega_{\mathfrak p}
\]
vaut pour toutes les algèbres \(A\) du groupe \(\mathfrak K\) associé à \(K\).

\paragraph{Preuve.}
Comme \(f\) est le \(\mathfrak p\)-degré de \(K\) (le même pour tous les diviseurs premiers de \(\mathfrak p\) dans \(K\)), tandis que, d'autre part, le \(\mathfrak p\)-indice de \(A\) est égal à l'indice, donc à l'exposant, de \(A_{\mathfrak p}\), il résulte du théorème 3 que \(f\) est en tout cas un multiple de ce plus petit exposant; il reste donc seulement, pour la preuve, à montrer qu'il existe dans le groupe \(\mathfrak K\) des algèbres \(A\) ayant exactement le \(\mathfrak p\)-indice \(f\).

D'après le théorème de densité de Frobenius, il existe sûrement encore un autre idéal premier \(\mathfrak p'\) de \(\Omega\), ne divisant pas le discriminant relatif de \(K\), dont les diviseurs premiers dans \(K\) ont le degré relatif \(f\). Soit alors \(Z\) un corps cyclique sur \(\Omega\) dont le \(\mathfrak p\)-degré et le \(\mathfrak p'\)-degré aient la valeur \(f\). Soit de plus \(\alpha\) un nombre de \(\Omega\), pour lequel les symboles de reste normique
\[
\left(\frac{\alpha,Z}{\mathfrak p}\right)\quad\text{et}\quad
\left(\frac{\alpha,Z}{\mathfrak p'}\right)
\]
aient des valeurs réciproques d'ordre \(f\), tandis que pour tous les autres diviseurs \(\mathfrak q\) du conducteur de \(Z\) on a
\[
\left(\frac{\alpha,Z}{\mathfrak q}\right)=1.
\]
D'après le théorème généralisé de la progression arithmétique, \(\alpha\) peut de plus être choisi de façon qu'en dehors éventuellement de \(\mathfrak p,\mathfrak p'\) et des \(\mathfrak q\), il ne contienne plus qu'un seul idéal premier \(\mathfrak r\) de \(\Omega\), exactement à la première puissance. Pour celui-ci, d'après le théorème du produit pour le symbole de reste normique (loi de réciprocité), on a également
\[
\left(\frac{\alpha,Z}{\mathfrak r}\right)=1
\]
[H, 3.8--10]. Toute algèbre
\[
        (\alpha,Z)=A
\]
a alors le \(\mathfrak p\)-indice et le \(\mathfrak p'\)-indice \(f\), mais pour toutes les autres places premières de \(\Omega\) l'indice \(1\) [H, 17.7]. D'après le théorème 3, il en résulte, en tenant compte de l'hypothèse sur \(\mathfrak p\) et du choix de \(\mathfrak p'\), que \(K\) est corps de décomposition pour \(A\). Ainsi sont effectivement mises en évidence dans le groupe \(\mathfrak K\) des algèbres \(A\) de \(\mathfrak p\)-indice \(f\).

\paragraph{Théorème 5.}
\emph{(Théorème d'unicité et d'ordre.)} Si \(K\leqq K'\), alors \(\mathfrak K\leqq\mathfrak K'\), et réciproquement.

En particulier, l'association des groupes d'algèbres \(\mathfrak K\) aux corps galoisiens \(K\) est donc biunivoque.

\paragraph{Preuve.}
a) De \(K\leqq K'\), il résulte trivialement que \(\mathfrak K\leqq\mathfrak K'\): toute algèbre \(A\) décomposée par \(K\) est, a fortiori, décomposée par \(K'\).

b) Soit réciproquement \(\mathfrak K\leqq\mathfrak K'\). Alors, d'après le théorème de décomposition, on a \(f_{\mathfrak p}\mid f'_{\mathfrak p}\) pour tout non-diviseur \(\mathfrak p\) du discriminant relatif. En particulier, on a donc \(f_{\mathfrak p}=1\) pour presque tous les \(\mathfrak p\) tels que \(f'_{\mathfrak p}=1\). Mais il en résulte, par le procédé analytique usuel\footnote{Voir à ce sujet par exemple H. Hasse [loc. cit. note 6], \S 25, III.}, que \(K\leqq K'\).

\paragraph{8.}
Enfin, le théorème principal apporte un progrès essentiel à la question, étudiée par I. Schur\footnote{I. Schur, \emph{Arithmetische Untersuchungen ueber endliche Gruppen linearer Substitutionen}, Berl. Akad.-Ber. 1906.}, des corps de nombres dans lesquels les représentations absolument irréductibles d'un groupe fini sont réalisables:

\paragraph{Théorème 6.}
Les représentations absolument irréductibles d'un groupe fini \(\mathfrak G\) sont toutes possibles dans des corps cyclotomiques, par exemple toujours dans le corps des racines \(n^h\)-ièmes de l'unité, si \(n\) est l'ordre de \(\mathfrak G\) et si \(h\) est suffisamment grand.

\paragraph{Preuve.}
Si l'on passe à l'algèbre de groupe \(G=\mathbb Q[\mathfrak G]\), les représentations absolument irréductibles \(\Gamma_i\) de \(\mathfrak G\) deviennent les représentations absolument irréductibles des constituants simples \(G_i\) de l'algèbre semi-simple \(G\), et leurs centres sont chaque fois les corps \(\Omega_i\) des caractères correspondants,\footnote{Voir à ce sujet: a) R. Brauer et E. Noether, \emph{Ueber minimale Zerfaellungskoerper irreduzibler Darstellungen}, Berl. Akad.-Ber. 1927; \S 1. b) R. Brauer, \emph{Ueber Systeme hyperkomplexer Zahlen}, Math. Zeitschr. 30 (1929); théorème 3. c) E. Noether, \emph{Hyperkomplexe Groessen und Darstellungstheorie}, Math. Zeitschr. 30 (1929); \S\S 21, 24, 26.} donc, en raison de la finitude de \(\mathfrak G\), des corps cyclotomiques, et même certainement des sous-corps du corps des racines \(n\)-ièmes de l'unité.

D'après le théorème 3, un corps cyclique \(Z_i\) sur \(\Omega_i\) est maintenant corps de décomposition pour \(G_i\), si, pour toute place première \(\mathfrak p\) de \(\Omega_i\), son \(\mathfrak p\)-degré \(n_{i\mathfrak p}\) est un multiple du \(\mathfrak p\)-indice \(m_{i\mathfrak p}\) de \(G_i\). \(m_{i\mathfrak p}\) n'est différent de 1 que pour les diviseurs premiers de l'idéal fondamental de \(G_i\) relativement à \(\Omega_i\), donc certainement seulement pour les diviseurs premiers du discriminant absolu de \(G_i\); comme celui-ci divise \(n^n\) -- \(n^n\) est le discriminant d'un ordre (non maximal) dans \(G\)\footnote{Voir E. Noether [loc. cit. note 12c], \S 26.} --, donc tout au plus pour les diviseurs premiers \(\mathfrak p\) de \(n\). Pour faire de \(n_{i\mathfrak p}\), pour ces \(\mathfrak p\), un multiple de \(m_{i\mathfrak p}\), il suffit de prescrire que les \(\mathfrak p\) en question soient ramifiés dans \(Z_i\) avec un ordre chaque fois divisible par \(m_{i\mathfrak p}\). Or c'est ce que réalise, comme on le voit facilement, le corps des racines \(n^h\)-ièmes de l'unité pour \(h\) suffisamment grand.

Dans la dernière phrase du travail cité, I. Schur constate que, dans tous les cas connus jusqu'ici, le corps des racines \(n\)-ièmes de l'unité suffit déjà. La question de savoir si cela vaut toujours, et si les méthodes développées ici suffisent pour trancher cette question, reste réservée à des recherches ultérieures.

\begin{center}
Reçu le 11 novembre 1931.
\end{center}

\setcounter{footnote}{0}
\clearpage
